\Needspace{8\baselineskip}
\section{综合算例与关键结论}
\label{l8:sec:examples}
\subsection{末尾不足一个完整块时的计算}
沿用一维输入和滤波器，逐项分析末尾块的载入与计算。取 $W=7$、$B=4$、$K=5$、$r=2$，输入 $N=[1,2,3,4,5,6,7]$，滤波器 $M=[3,4,5,4,3]$。网格需要两块，分别负责输出下标 $0$—$3$ 与 $4$—$7$，其中 $P[7]$ 没有对应的有效输出位置。

第二块从 $s=4$ 开始。策略 1 的共享数组覆盖全局位置 $2$ 到 $9$，依次填入 $[3,4,5,6,7,0,0,0]$。前五个值来自真实输入，后三个值按零填充规则生成。三段载入中，线程 $t=3$ 虽然对应无效输出 $i=7$，仍负责把 $N[3]=4$ 放入共享位置 $1$，并把其核心位置写为零。

如果线程 $t=3$ 在搬运前提前返回，共享位置 $1$ 就没有被它初始化。有效线程 $t=0$ 计算 $P[4]$ 时需要共享位置 $0$ 到 $4$，其中就包含这一位置。因此有效输出会依赖未完成的搬运；只查看“哪些线程要写输出”无法判断哪些线程可以提前结束。

\begin{table}[H]
\centering\small
\begin{tabularx}{\textwidth}{@{}p{.12\textwidth}p{.15\textwidth}p{.34\textwidth}Y@{}}
\toprule
局部线程 & 输出位置 & 策略 1 读取的共享值 & 结果与写入 \\
\midrule
$t=0$ & $P[4]$ & $[3,4,5,6,7]$ & $9+16+25+24+21=95$ \\
$t=1$ & $P[5]$ & $[4,5,6,7,0]$ & $12+20+30+28=90$ \\
$t=2$ & $P[6]$ & $[5,6,7,0,0]$ & $15+24+35=74$ \\
$t=3$ & 无有效输出 & $[6,7,0,0,0]$ & 可执行乘加，但不写 $P[7]$ \\
\bottomrule
\end{tabularx}
\caption{末尾块的完整窗口；零值也占据明确的共享位置。}
\end{table}

两轮载入写法中，第二块第一轮载入 $[N[2],N[3],N[4],N[5]]$，第二轮载入 $[N[6],0,0,0]$，最终共享内容与三段写法相同。在不同写法中负责同一数据的线程可能不同，但只要式\eqref{l8:eq:sharedmapping}成立，乘加结果就一致。

对策略 3，同一块的共享内容只有 $[N[4],N[5],N[6],0]=[5,6,7,0]$。计算 $P[4]$ 时从全局地址取得 $N[2],N[3]$，从共享数组取得 $N[4],N[5],N[6]$。不同数据放置方式由不同下标转换衔接到同一个卷积窗口。

\subsection{带边界条件的精确计数}
对直接核函数，位置 $i$ 的有效乘加项数就是窗口 $[i-r,i+r]$ 与数组 $[0,W-1]$ 的交集长度。将边界判断写成计数式，可得
\begin{equation}
 v_i=\min(W-1,i+r)-\max(0,i-r)+1,
 \qquad 0\le i<W.
\label{l8:eq:validcount}
\end{equation}
对于 $W=7,r=2$，$v_i$ 依次为 $3,4,5,5,5,4,3$，总计 $29$。直接核函数每个有效项读一个输入和一个系数，因此一共执行 $58$ 次标量读取，对应 $232$ B 的逻辑读入；有效输出写入 $7$ 个 \mbox{\code{float}}，对应 $28$ B；乘法加法合计为 $58$ FLOPs。这些数字来自代码实际通过边界分支的项数。

若只比较策略 1 对输入 $N$ 的请求，首块的输入覆盖 $[-2,5]$，其中六个位置有效；第二块覆盖 $[2,9]$，其中五个位置有效，共读入 $11$ 个真实输入值，即 $44$ B。共享内存中总共写入 $2\times8=16$ 个位置，其余五个位置由线程直接写零。常量系数的读取和缓存事务没有包含在这 $44$ B 中。

完整输入块已在载入阶段填零，策略 1 的乘加循环就不再跳过零项。七个有效输出各执行五次乘加，需要 $70$ FLOPs；代码\ref{l8:lst:three}与代码\ref{l8:lst:two}采用全部线程计算的结构，八个已启动线程共执行 $80$ FLOPs，最后一个线程的结果不写出。\textbf{减少输入请求可以同时增加零项或无效输出上的算术工作}，这正是按计数对象分别核对的意义。

再考虑 $W=16,B=4,K=5$ 的布局。四个块的有效输入请求数依次为 $6,8,8,6$，策略 1 合计请求 $28$ 个输入值。直接核函数的有效输入使用次数为 $3+4+12\times5+4+3=74$。该例按整个网格统计的输入请求减少因子为 $74/28\approx2.64$，而式\eqref{l8:eq:reuse}的 $2.5$ 比较的是一个远离边界的完整内部块。两者采用了不同的统计范围。

\Needspace{24\baselineskip}
\subsection{关键公式与实现关系}
卷积计算、缓存与分块可以通过以下关系一起核对。公式中的 $K=2r+1$，$B$ 表示每块输出数，$s=bB$ 表示输出块起点；策略 2 若按输入块大小配置线程，线程数应相应增加为 $B+2r$。

\begin{table}[H]
\centering\small
\begin{tabularx}{\textwidth}{@{}p{.30\textwidth}Y@{}}
\toprule
对象 & 关系 \\
\midrule
输出与输入窗口 & $P[i]=\sum_{j=0}^{K-1}\widetilde N[i-r+j]M[j]$ \\
一块所需输入数 & $T_{\mathrm{in}}=B+2r=B+K-1$ \\
完整输入块的共享映射 & $N_{\mathrm{ds}}[q]=\widetilde N[s-r+q]$ \\
策略 1 的计算下标 & $q=t+j$ \\
核心区的共享映射 & $N_{\mathrm{ds}}[q]=\widetilde N[s+q]$ \\
策略 3 的共享计算下标 & $q=t-r+j$，仅当 $s\le i-r+j<s+B$ \\
内部块输入请求比较 & 直接法 $BK$；策略 1 为 $B+2r$ \\
策略 1 的线程覆盖条件 & 三段载入要求 $r\le B$；两轮载入要求 $2r\le B$ \\
共享内存容量 & 策略 1、2 为 $4(B+2r)$ B；策略 3 为 $4B$ B \\
\bottomrule
\end{tabularx}
\caption{卷积窗口、共享索引与资源用量。}
\end{table}

从输出窗口出发确定完整输入范围，再选择共享数组的起点，即可推得载入与计算下标。常量存储器处理统一系数的读取，共享内存处理邻域输入的复用；二者连接到同一个卷积计算，而优化的对象各不相同。
\FloatBarrier
